\section{Evaluation Metrics}
\label{sec:eval}
In recent work, video generation models have largely been assessed using image-derived metrics such as Inception Score (IS) \cite{Salimans2016}, Fréchet Inception Distance (FID) \cite{Heusel2017} and its temporal extension Fréchet Video Distance (FVD) \cite{Unterthiner2019}, along with Structural Similarity Index (SSIM) \cite{wang2004image} and Learned Perceptual Image Patch Similarity (LPIPS) \cite{zhang2018unreasonable}, as well as text alignment via CLIPScore \cite{Hessel2021}. While these measures offer convenient benchmarks, they obscure crucial factors such as temporal coherence, storytelling fidelity and multi-scene consistency, and frequently diverge from human judgments on longer, more complex clips.

To address these shortcomings, VBench introduces a comprehensive, hierarchical evaluation suite that decomposes “video generation quality” into fine-grained dimensions such as visual quality, motion smoothness, identity consistency, temporal flicker, spatial relationships and text-video relevance. Each of these dimensions are driven by tailored prompt sets and validated with human preference annotations \cite{Huang2024}. By providing multi-dimensional scores rather than a monolithic metric, VBench enables detailed diagnostics of generative strengths and weaknesses, making it the principal benchmark for next-generation video models. Another similar benchmark that is multi-dimensional is Wan2.1 \cite{wan2025}, however, VBench has been used by most of the recent literature.